\section{Competencia, especialización y diversidad}

\subsection{Estrecho y fuerte vs. amplio y superficial}

\begin{importantbox}{La superinteligencia puede ser estrecha}
Ilya propuso una posibilidad que suele pasarse por alto: la superinteligencia no tiene que ser un "oráculo que sabe hacer todo"; puede ser \textbf{useful and narrow at the same time}. En el futuro podría haber muchas superinteligencias estrechas distintas, cada una especializada en un dominio diferente. La competencia se desarrollará mediante la especialización (specialization), tal como ocurre en los mercados y en la evolución.
\end{importantbox}

\subsection{El problema de la diversidad en la IA}

La razón por la que la IA actual carece de diversidad está en el pre-training: todos los modelos se entrenan con los mismos datos, por lo que sus resultados son sorprendentemente parecidos entre sí. El RL y el post-training empiezan a introducir cierta diferenciación.

\begin{knowledgebox}{Self-play y diversidad adversarial}
El self-play atrajo el interés de Ilya en un principio porque ofrecía una manera de generar señal de entrenamiento "solo con compute, sin data". Pero el self-play clásico solo era bueno para desarrollar habilidades específicas (negociación, conflicto, estrategia). Hoy, el self-play regresa en formas distintas: debate, verifier, LLM-as-judge y otros esquemas adversariales.

La competencia genera diversidad de manera natural: cuando un agent ve que otro agent ya adoptó cierto método, tiene un incentivo para explorar direcciones diferentes.
\end{knowledgebox}

\subsection{Predicción de timeline}

Ilya dio un rango de \textbf{5 a 20 años} para llegar a un agent con aprendizaje continuo similar al humano (y, a partir de ahí, a la superinteligencia).

\subsection{Resumen de la sección}

La superinteligencia puede ser estrecha y fuerte a la vez. La competencia impulsará la especialización y la diversidad. El pre-training explica la homogeneidad de los modelos actuales, y el RL empieza a introducir diferencias. Ilya prevé que la IA alcance una capacidad de aprendizaje similar a la humana en un plazo de 5 a 20 años.

